\documentclass{caosp306}
\usepackage{graphicx}
\usepackage{natbib}
\bibliographystyle{caosp306}

%%%%%%%%%%%%%%%%%%%%%%%%%%%%%%%%%%%%%%%%%%%%%%%%%%%%%%%%%%%%%%%%%%%%%%%%%%%%%
%                       E D I T O R I A L   N O T E S                        
% Next 7 lines are used by the editors, who will fill-in proper values. 
%%%%%%%%%%%%%%%%%%%%%%%%%%%%%%%%%%%%%%%%%%%%%%%%%%%%%%%%%%%%%%%%%%%%%%%%%%%%%
\articleNo{AXRO01}
\pubyear{2018}
\volume{49}
\volnumber{2}
\firstpage{1}
\received{}
\accepted{}   


\def\BibTeX{{\rm B\kern-.05em{\sc i\kern-.025em b}\kern-.08em
             T\kern-.1667em\lower.7ex\hbox{E}\kern-.125emX}}

%%%%%%%%%%%%%%%%%%%%%%%%%%%%%%%%%%%%%%%%%%%%%%%%%%%%%%%%%%%%%%%%%%%%%%%%%%%%%
\begin{document}

%%%%%%%%%%%%%%%%%%%%%%%%%%%%%%%%%%%%%%%%%%%%%%%%%%%%%%%%%%%%%%%%%%%%%%%%%%%%%
%              R U N N I N G   P A G E   H E A D I N G S                     
% Odd page headings (except for the title page) are produced automatically
% and contain the title. If, and only if, the title of your article is too
% long the running head is ommitted in the printout; you can make your own
% running title by using the \htitle command, putting the shortened title
% between the curly brackets.
%
% Unlike the \htitle command, the \hauthor command is compulsory. It is
% used to produce even page headings and contains the names of the authors
% of an article.  All authors must be listed here, if possible. When
% authors' list is too long, you can abbreviate it by using "{\it et
% al.}". Authors' names are given in the form: initial(s) of the author's
% first name and surname. Authors are separated by a "," (comma) sign and
% the last one by "and".
%%%%%%%%%%%%%%%%%%%%%%%%%%%%%%%%%%%%%%%%%%%%%%%%%%%%%%%%%%%%%%%%%%%%%%%%%%%%%
%\htitle{An example of shorter title ...}

\hauthor{F.\,Irst, S.\,Econd, and T.\,Hird}
%\hauthor{F.\,Irst {\it et al.}}
%%%%%%%%%%%%%%%%%%%%%%%%%%%%%%%%%%%%%%%%%%%%%%%%%%%%%%%%%%%%%%%%%%%%%%%%%%%%%
%                       T I T L E                                            
% Capital letters in the title are only used at the beginning of the
% names. Don`t end the title by a "." (dot)
%%%%%%%%%%%%%%%%%%%%%%%%%%%%%%%%%%%%%%%%%%%%%%%%%%%%%%%%%%%%%%%%%%%%%%%%%%%%%
\title{An example of title}

%%%%%%%%%%%%%%%%%%%%%%%%%%%%%%%%%%%%%%%%%%%%%%%%%%%%%%%%%%%%%%%%%%%%%%%%%%%%%
%                   A U T H O R  N A M E S                                   
% Authors' names are separated by the \and command and their institutes
% are assigned by the \inst{n} command.

%%%%%%%%%%%%%%%%%%%%%%%%%%%%%%%%%%%%%%%%%%%%%%%%%%%%%%%%%%%%%%%%%%%%%%%%%%%%%
\author{
        F.\,Irst \inst{1,} \inst{2} 
      \and 
        S.\,Econd \inst{1}   
      \and 
        T.\,Hird \inst{3}
       }

%%%%%%%%%%%%%%%%%%%%%%%%%%%%%%%%%%%%%%%%%%%%%%%%%%%%%%%%%%%%%%%%%%%%%%%%%%%%%
%           I N S T I T U T E S'  A D D R E S S E S                          
% The affiliation of authors is generated by the \institute command, the
% \and command being again used to separate individual addresses.
% The given postal address must be complete in order to facilitate our
% editorial work. Moreover, you can add your e-mail address, using the
% \email command.
%%%%%%%%%%%%%%%%%%%%%%%%%%%%%%%%%%%%%%%%%%%%%%%%%%%%%%%%%%%%%%%%%%%%%%%%%%%%%
\institute{
          Czech Technical University in Prague,\\
         Technicka 2, Prague 166 27, Czech Republic \email{sample@cvut.cz}
         \and
           \ondrejov
         \and
         	Aschaffenburg University of Applied Sciences,\\
            Wuerzburger Strasse 45, 63743 Aschaffenburg, Germany
          }

%%%%%%%%%%%%%%%%%%%%%%%%%%%%%%%%%%%%%%%%%%%%%%%%%%%%%%%%%%%%%%%%%%%%%%%%%%%%%
%                        D A T E / R E C E I V E D                          
% Date inserted here will be the date when your paper was received The
% format is: month (not abbreviated), day, year.
%%%%%%%%%%%%%%%%%%%%%%%%%%%%%%%%%%%%%%%%%%%%%%%%%%%%%%%%%%%%%%%%%%%%%%%%%%%%%
\date{April 7, 2017}

%%%%%%%%%%%%%%%%%%%%%%%%%%%%%%%%%%%%%%%%%%%%%%%%%%%%%%%%%%%%%%%%%%%%%%%%%%%%%
%                        M A K E T I T L E
% The beginning part (title, author(s), etc.) of your article must be
% closed by the \maketitle command.
%%%%%%%%%%%%%%%%%%%%%%%%%%%%%%%%%%%%%%%%%%%%%%%%%%%%%%%%%%%%%%%%%%%%%%%%%%%%%
\maketitle

%%%%%%%%%%%%%%%%%%%%%%%%%%%%%%%%%%%%%%%%%%%%%%%%%%%%%%%%%%%%%%%%%%%%%%%%%%%%%
%                        A B S T R A C T,  K E Y W O R D S                   
% Keywords must be separated from each other by -- (two dashes).
%%%%%%%%%%%%%%%%%%%%%%%%%%%%%%%%%%%%%%%%%%%%%%%%%%%%%%%%%%%%%%%%%%%%%%%%%%%%%
\begin{abstract}
This template can be used to prepare a proceeding for the 10\textsuperscript{th} International Workshop on Astronomical X--Ray Optics Prague
\keywords{Keyword 1 -- Keyword 2 -- Keyword 3}
\end{abstract}

%%%%%%%%%%%%%%%%%%%%%%%%%%%%%%%%%%%%%%%%%%%%%%%%%%%%%%%%%%%%%%%%%%%%%%%%%%%%%
%                       S E C T I O N I N G                                  
%
% The (sub)section titles are typeset in boldface; so, if working in the
% mathematics mode in (sub)section titles, you must use \boldmath and 
% enclose it into curly brackets, e.g. "{\bolmath $R^{2}$}".
%%%%%%%%%%%%%%%%%%%%%%%%%%%%%%%%%%%%%%%%%%%%%%%%%%%%%%%%%%%%%%%%%%%%%%%%%%%%%
\section{Examples of Article Components}
%%%%%%%%%%%%%%%%%%%%%%%%%%%%%%%%%%%%%%%%%%%%%%%%%%%%%%%%%%%%%%%%%%%%%%%%%%%%%
%                       L A B E L                                            
% The label command is very convenient for you when referring to sections,
% subsections,..., tables, figures as well as to equations (see commands
% \ref and \pageref). In the case of figure and/or table environments the
% \label command should always be put after the \caption command to
% preserve proper numbering. When using the \label command the file must
% be compiled twice to get proper cross-references.
%%%%%%%%%%%%%%%%%%%%%%%%%%%%%%%%%%%%%%%%%%%%%%%%%%%%%%%%%%%%%%%%%%%%%%%%%%%%%
\label{sec:examples}

The sections below show examples of different article components.
%%%%%%%%%%%%%%%%%%%%%%%%%%%%%%%%%%%%%%%%%%%%%%%%%%%%%%%%%%%%%%%%%%%%%%%%%%%%%
%                       P A R A G R A P H                                    
% To generate a paragraph simply leave a blank line after the last
% sentence of the preceding paragraph as shown below.
%%%%%%%%%%%%%%%%%%%%%%%%%%%%%%%%%%%%%%%%%%%%%%%%%%%%%%%%%%%%%%%%%%%%%%%%%%%%%

As internal references use label commands \verb|\label{}| for section, subsection, talbes and figures. This command used in pair with \verb|\ref{}| command. 

An example: "...see section~\ref{sec:examples}."

Always remember the following rules:
\begin{itemize}
  \item physical and mathematical variables are typeset in Italics
  \item physical units and mathematical functions are typeset in Roman
  \item instead using "km/s" use "km\,s$^{-1}$"
  \item use small spaces: not "150 AU", but "150\,AU". Moreover, this small space cannot be line-breaked.
\end{itemize}

Note the following examples:
 inside normal text: {\it v} = 5\,km\,s$^{-1}$ 
 inside mathematics mode/environment: $ v = 5\,{\rm km\,s}^{-1}$ 

%%%%%%%%%%%%%%%%%%%%%%%%%%%%%%%%%%%%%%%%%%%%%%%%%%%%%%%%%%%%%%%%%%%%%%%%%%%%%
%                       S U B S E C T I O N S                                
% begin with the \subsection command; the remaining rules are identical
% with those introduced for sections.
%%%%%%%%%%%%%%%%%%%%%%%%%%%%%%%%%%%%%%%%%%%%%%%%%%%%%%%%%%%%%%%%%%%%%%%%%%%%%

\subsection{Tables}
%%%%%%%%%%%%%%%%%%%%%%%%%%%%%%%%%%%%%%%%%%%%%%%%%%%%%%%%%%%%%%%%%%%%%%%%%%%%%
%                       T A B L E S                                          

% Inside the tabular enviroinment try to avoid commands like \noindent,
% \hspace.
%
% A table caption must be placed above the actual table (more details
% about the caption are given bellow).
%
% As to a table layout, use the following rules:
% - Do not use vertical lines at table edges. If needed they can be used
%   between columns.
% - A table must always start and end with two horizontal lines
%   (\hline\hline). This rule applies even to tables splitted into several
%   pages: two \hline are used just twice (the beginning and the end), not
%   at the bottom of the table when it continues on the following page.
% - Do not separate every table row by a horizontal line, use the latter
%   only when necessary.
%%%%%%%%%%%%%%%%%%%%%%%%%%%%%%%%%%%%%%%%%%%%%%%%%%%%%%%%%%%%%%%%%%%%%%%%%%%%%
The caption of a table is created by the \verb|\caption| command. Do not number
it, nor write Tab. or Table - let again \LaTeX~do it for you.  The
\verb|\caption| command should always be used before the table body and should
not exceed two lines in length. The caption should end with a "." (dot).

\begin{table}[htbp]
\begin{center}
	\caption{A sample of table.}
	\label{tab:t1}
	\begin{tabular}{llll}
		\hline\hline
		planet  &{\it u}&    {\it h}      & $p_{\uparrow}M_{\rm K}$\\
		\hline
		Jupiter & 0.08  & 0.03 $\pm$ 0.01 & 2.75$nM_{\rm KP}$ \\
		Saturn  & 0.16  & 0.16 $\pm$ 0.04 & 1.16$nM_{\rm KP}$ \\
		Uranus  & 0.24  & 1.3  $\pm$ 0.5  & 0.43$nM_{\rm KP}$ \\
		Neptune & 0.52  & 2.6  $\pm$ 0.7  & 0.72$nM_{\rm KP}$ \\
		\hline\hline
	\end{tabular}
\end{center}
\end{table}

Tables in the text should be referred as "Tab.", e.g. "... Tab.\,\ref{tab:t1} shows ...". 

Because figures and tables are "floating bodies" their position in the text may not
always be proper - the style file tries to put it, whenever possible, on
top of the page. If you wish to put a table at a particular place in the
text you can use the options [b], [h], [t]. It is acknowledged that
positioning of tables is by no means straightforward.

\subsection{Figures}   
%%%%%%%%%%%%%%%%%%%%%%%%%%%%%%%%%%%%%%%%%%%%%%%%%%%%%%%%%%%%%%%%%%%%%%%%%%%%%
%                       F I G U R E S     
Figures must be stored in an EPS (Encapsulated PostScript).
Authors are required to produce their EPS graphics in a vector mode a
bitmap version can be used only if absolutely necessary, e.g. for
photos. Using this way of preparing figures you are kindly asked to submit the
figures in EPS files to the editors, along with the \LaTeX~ source file of
your article. 

Colour figures are forbidden.

Maximum dimensions of a single picture are horizontal length: 11.5\,cm and vertical length: 17.0\,cm.                                    
The size of the axes' description should not be less than 0.3\,cm (alpha-
numeric signs). Ticks on axes must lie in the inner side of axes. The resulting line width should be 0.1 - 0.3\,mm.

Figures can be inserted into the text by the command \verb|\includegraphics|
(see graphicx package documentation for available options).

The \verb|\caption| command should always be used after the figure body.

\begin{figure}[htbp]
\centering
	\includegraphics[width=\linewidth]{example-image}
	\caption{This is a picture from graphic file formatted in EPS format loaded by the command ~$\backslash$includegraphics}
	\label{fig:1}
\end{figure}

A more complicated example of a figure consisting of 2 separate graphics files is given in Fig.\,\ref{fig:2}.

\begin{figure}[htbp]
\centerline{\includegraphics[width=5.cm]{example-image-a} 
           \hspace{0.5cm}
           \includegraphics[width=5.cm]{example-image-b}
          }
\caption{This is a more complicated example of a figure.}
\label{fig:2}
\end{figure}

All figures and tables should be referenced in the text and should be placed on the page where they are first cited or in subsequent pages.

\subsection{Sample Equation}                   
Now, let us introduce a few examples of how to write equations; for
further (technical) details, please, consult any \LaTeX~ manual(s).

If the equation is splitted into several lines, always number it on its
last line. In this case, you must take care to allow the proper size of
parentheses, brackets, etc. when the corresponding pair is not placed on
the same line: use the \verb|\vphantom| command.

In the case of the array of equations the number goes on its middle
line. In case of a two-member array use the second line for its
numbering.

\begin{equation}
 \label{eq:1}
E_{\rm 1}=-GmM_{\odot}/(2a_{\rm o}),
\end{equation}
where $G$ is the gravitational constant, $m$ is the mass of comet,
$M_{\odot}$ - the mass of the Sun. 
References will work as you'd expect: (\ref{eq:1}).



\subsection{Citations}
To cite an article, book, proceedings, please use generally adopted
rules, e.g. \citep{2005ApJ...623..411B}. If there are more citations inside the parenthesis, separate
them by a "," (comma): \citep{2005ApJ...623..411B, 1687pnpm.book.....N}. 

Next line show the examples of basic \BibTeX~ usage in citations:

\begin{itemize}
  \item \verb|\citep{key}| shows \citep{2005ApJ...623..411B}
\end{itemize}


\section{References}
References must be sorted out in \textbf{alphabetical order}.
\\
Irrespective of the number of authors you must list each of them. Within
this environment it is not allowed to use "{\it et al.}".
\\
Create a database (.bib) file that describes the articles or books you 
   want to reference. The NASA Astrophysics Data System (ADS) provides
   automatic tools for retrieving a .bib file including entries for a
   selection of articles.  An example of a typical .bib file is also 
   provided in the CAOSP \LaTeX~ macro package.
   If there is no entry in ADS for some of your references, you are
   obliged to add these entries into your .bib file "by hand" respecting
   the formal structure of the record. Standard journal abbreviations
   defined by ADS are required.
   \\
   To ensure that each reference is spelled properly, and to facilitate ADS
crossreferencing as well, every author is also obliged to supply an
ADS-identification entry of each reference. This ADS-identification entry
has the form of a \LaTeX~ comment just in front of a given reference on a
separate line, which has to start with a key-word "ADS\_ID". 
%In order to find a specific ADS-identification, please use the following on-line form http://adsabs.harvard.edu/cgi-bin/nph-abs_connect
% An ADS-identification entry is a unique string whose length is exactly 19
% characters, e.g. 1955ApJ...121..349B

% You should make use of the journal abbreviations listed in alphabetical
% order below.
%  command     abbreviation                             %journal name                               
%  ~~~~~~~     ~~~~~~~~~~~~                             ~~~~~~~~~~~~~                               
%                                                                            
% \abasb	Abastumani Bull.						% Abastumani Bulletin
% \aaa		Astron. Astrophys.						% Astronomy and Astrophysics
% \aaas		Astron. Astrophys., Suppl. Ser.			% Astronomy and Astrophysics, Supplement Series
% \aap		Astron. Astrophys.						% Astronomy & Astrophysics
% \aapr		Astron. Astrophys. Rev.					% Astronomy & Astrophysics Reviews
% \aaps		Astron. Astrophys., Suppl.				% Astronomy & Astrophysics, Supplement
% \aas		Acta Astron. Sinica						% Acta Astronomica Sinica
% \aca		Acta Astron.							% Acta Astronomica
% \aipconf	Amer. Inst. Phys. Conf. Proc.			% American Institute of Physics Conference Proceedings
% \aj		Astron. J.								% Astronomical Journal
% \an		Astron. Nachr.							% Astronomische Nachrichten
% \ao		Appl. Opt. 								% Applied Optics
% \apj		Astrophys. J.							% Astrophysical Journal
% \apjl		Astrophys. J., Lett.					% Astrophysical Journal, Letters
% \apjlett	Astrophys. J., Lett.					% The Astrophysical Journal, Letters
% \apjs		Astrophys. J., Suppl.					% The Astrophysical Journal, Supplement
% \apjss	Astrophys. J., Suppl. Ser.				% Astrophysical Journal, Supplement Series
% \apjsupp	Astrophys. J., Suppl.					% The Astrophysical Journal, Supplement
% \aplett	Astrophys. Lett.						% Astrophysics Letters
% \applopt	Appl. Opt.								% Applied Optics
% \apspr	Astrophys. Space Phys. Res.				% Astrophysics Space Physics Research
% \apss		Astrophys. Space Sci.					% Astrophysics and Space Science
% \ar		Ann. Rev. Astron. Astrophys.			% Annual Review of Astronomy and Astrophysics
% \araa		Ann. Rev. Astron. Astrophys.			% Annual Review of Astronomy & Astrophysics
% \aspconf	ASP Conf. Ser.							% Astronomical Society of the Pacific Conference Series
% \asr		Adv. Space Res.							% Advances in Space Research
% \asslconf	Astrophys. Space Sien. Lib. Conf. Ser.	% Astrophysics & Space Science Library Conference Series
% \astap	Astron. Astrophys.						% Astronomy & Astrophysics
% \azh		Astron. Zh.								% Astronomicekij Zhurnal
% \baas		Bull. Am. Astron. Soc.					% Bulletin of the American Astronomical Society
% \bac		Bull. Astron. Inst. Czechosl.			% Bulletin of the Astronomical Institutes of Czechoslovakia
% \bain		Bull. Astron. Inst. Netherlands			% Bulletin of the Astronical Instutes of Netherlands
% \basi		Bull. Astron. Soc. India				% Bulletin of the Astronical Society of India
% \caa		Chinese Astron. Astrophys.				% Chinese Astronomy & Astrophysics
% \caos		Contrib. Astron. Obs. Skalnat\'e Pleso	% Contributions of the Astronomical Observatory Skalnate Pleso
% \cpuo		Contrib. Princeton Univ. Obs.			% Contributions Princeton University Observatory
% \emp		Earth, Moon, Planets					% Earth, Moon, Planets
% \fcp		Fund. Cosmic Phys.						% Fundamental Cosmic Physics
% \harva	Harvard Ann.							% Harvard Annales
% \harvc	Harvard Circ.							% Harvard Circular
% \iauc		IAU Circ.								% IAU Circular
% \iaucirc	IAU Circ.								% IAU Circulars
% \iaut		IAU Trans.	
% \ibvs		Inf. Bull. Variable Stars				% Information Bulletin on Variable Stars
% \ica		Icarus									% Icarus
% \jrasc	J. R. Astron. Soc. Can.					% Journal of the Royal Astronomical Society of Canada
% \jrcan	J. R. Astron. Soc. Can.					% Journal of the Royal Astronomical Society of Canada
% \krym		Izv. Krymskoj Astrofiz. Obs.			% Izvestya Krymskoj Astrofizicheskoj Observatorii
% \lickb	Lick Bull.								% Lick Bulletin
% \maps		Met. Planet. Sci.						% Meteoritics & Planetary Science
% \memras	Mem. R. Astron. Soc. 					% Memoirs of the Royal Astronomical Society
% \mnras	Mon. Not. R. Astron. Soc.				% Monthly Notices of the Royal Astronomical Society
% \mpci		Minor Planet Circ.						% Minor Planet Circulars
% \mras		Mem. Royal Astron. Soc.					% Memoirs Royal Astron. Society
% \msai		Mem. Soc. Astron. Ital.					% Memoirs della Societa Astronomica Italiana
% \mvs		Mitt. Verand. Sterne					% Mitteilungen Verand. Sterne
% \nat		Nature									% Nature
% \obs		Observatory								% The Observatory
% \odesai	Odessa Izv.								% Odessa Izvestija
% \paics	Publ. Astron. Inst. Czechosl. Acad. Sci.% Publications of the Astronomical Institute of the Czechoslovak Academy of Sciences
% \pasj		Publ. Astron. Soc. Jap.					% Publications of the Astronomical Society of Japan
% \pasp		Publ. Astron. Soc. Pac.					% Publ. of the Astron. Society of the Pacific
% \pdao		Publ. DAO								% Publications of DAO
% \perzve	Perem. Zvezdy							% Peremennye Zvezdy
% \planss	Planet. Space Sci.						% Planetary Space Science
% \pra		Phys. Rev. A							% Physical Review A: Atomic, Molecular, and Optical Physics
% \prb		Phys. Rev. B							% Physical Review B: Condensed Matter and Materials Physics
% \prc		Phys. Rev. C							% Physical Review C: Nuclear Physics
% \prd		Phys. Rev. D							% Physical Review D: Particles, Fields, Gravitation and Cosmology
% \pre		Phys. Rev. E							% Physical Review E: Statistical, Nonlinear, and Soft Matter Physics
% \prl		Phys. Rev. Lett.						% Physical Review Letters
% \procspie	Proc. SPIE								% Proceedings of the International Society for Optical Engineering
% \pss		Planet. Space Sci.						% Planetary and Space Science
% \qjras	Quart. Jrn. Royal Astron. Soc.			% Quarterly Journal of the Royal Astronomical Society
% \rmaa		Rev. Mex. Astron. Astrofis.				% Revista Mexicana de Astronomia y Astrofisica.
% \rmp		Rev. Mod. Phys.							% Reviews of Modern Physics
% \sal		Sov. Astron. Lett.						% Soviet Astronomy Letters
% \sci		Science									% Science
% \sdan		Sol. dannye								% Solnechnye dannye
% \skytel	Sky Tel.								% Sky & Telescope
% \solphys	Sol. Phys. 								% Solar Physics
% \sonnv	Sonnenberg Veroff.						% Sonnenberg Veroff.
% \sph		Sol. Phys.								% Solar Physics
% \ssr		Space Sci. Rev.							% Space Science Reviews
% \va		Vistas Astron.							% Vistas in Astronomy

%                                                                            
% If you have not found the title of the journal among these commands you
% must type it yourself. In this case you are kindly asked to use the
% journal abbreviations as given in the ASTRONOMY and ASTROPHYSICS
% ABSTRACTS, Chapter " Periodicals, Proceedings, Books, Activities ".


%%%%%%%%%%%%%%%%%%%%%%%%%%%%%%%%%%%%%%%%%%%%%%%%%%%%%%%%%%%%%%%%%%%%%%%%%%%%%
%                       A C K N O W L E D G E M E N T S                      
% Next lines show you how to write acknowledgements.
% You must leave a blank line before the \acknowledgements command!
%%%%%%%%%%%%%%%%%%%%%%%%%%%%%%%%%%%%%%%%%%%%%%%%%%%%%%%%%%%%%%%%%%%%%%%%%%%%%

\acknowledgements
% Do not leave a blank line here! <---------------------->
This is the most common positions for acknowledgements.

%%%%%%%%%%%%%%%%%%%%%%%%%%%%%%%%%%%%%%%%%%%%%%%%%%%%%%%%%%%%%%%%%%%%%%%%%%%%%

\bibliography{example}


%%%%%%%%%%%%%%%%%%%%%%%%%%%%%%%%%%%%%%%%%%%%%%%%%%%%%%%%%%%%%%%%%%%%%%%%%%%%%
%                           A P P E N D I X                                  
%                       S P E C I A L   S I G N S                            
% Next table is to illustrate some special signs available in the CAOSP
% style.
%%%%%%%%%%%%%%%%%%%%%%%%%%%%%%%%%%%%%%%%%%%%%%%%%%%%%%%%%%%%%%%%%%%%%%%%%%%%%
\clearpage
\appendix{Special Signs}
You may wish to use special signs. A large number of the latter are listed
in the {\em {\LaTeX} User's Guide \& Reference Manual\/} by Leslie Lamport,
pp.\,44\,ff. We have created further symbols for math mode, and some special
Slovak characters (used e.g. in names), which cannot be simply produced by
{\LaTeX}.
\begin{table}[htbp]
\caption{Special signs}
\label{tspsig}
\begin{center}
\footnotesize
\renewcommand{\arraystretch}{1.25}
\begin{tabular}{lll@{\hspace{0.5cm}}lll}
\hline\hline
Input & Explanation & Output & Input & Explanation & Output\\
\hline
\verb|\la|     & less or approx       & $\la$     &
  \verb|\ga|     & greater or approx    & $\ga$\\
\verb|\getsto| & gets over to         & $\getsto$ &
  \verb|\cor|    & corresponds to       & $\cor$\\
\verb|\lid|    & less or equal        & $\lid$    &
  \verb|\gid|    & greater or equal     & $\gid$\\
\verb|\sol|    & similar over less    & $\sol$    &
  \verb|\sog|    & similar over greater & $\sog$\\
\verb|\lse|    & less over simeq      & $\lse$    &
  \verb|\gse|    & greater over simeq   & $\gse$\\
\verb|\grole|  & greater over less    & $\grole$  &
  \verb|\leogr|  & less over greater    & $\leogr$\\
\verb|\loa|    & less over approx     & $\loa$    &
  \verb|\goa|    & greater over approx  & $\goa$\\
\hline
\verb|\sun|      & sun symbol            & $\sun$     &
  \verb|\degr|     & degree                & $\degr$\\
\verb|\diameter| & diameter              & \diameter  &
  \verb|\sq|       & square                & \squareforqed\\
\verb|\fd|       & fraction of day       & \fd        &
  \verb|\fh|       & fraction of hour      & \fh\\
\verb|\fm|       & fraction of minute    & \fm        &
  \verb|\fs|       & fraction of second    & \fs\\
\verb|\fdg|      & fraction of degree    & \fdg       &
  \verb|\fp|       & fraction of period    & \fp\\
\verb|\farcs|    & fraction of arcsecond & \farcs     &
  \verb|\farcm|    & fraction of arcmin    & \farcm\\
\verb|\arcsec|   & arcsecond             & \arcsec    &
  \verb|\arcmin|   & arcminute             & \arcmin\\
\verb|\angstrom| & angstr\"{o}m          & \angstrom    &
                   &                       &       \\
\hline
\verb|\softL|    & Slovak L with caron   & \softL     &
  \verb|\softl|    & Slovak l with caron   & \softl\\
\verb|\softd|    & Slovak d with caron   & \softd     &
  \verb|\softt|    & Slovak t with caron   & \softt\\
\hline\hline
\end{tabular}
\renewcommand{\arraystretch}{1}
\end{center}
\end{table}


\end{document}
